%% International Journal of Digital Curation (IJDC) article starter, using
%% the official ijdc-v14 class from the dccpaper bundle.
%%
%% VENDORED: ijdc-v14.cls and dccpaper-base.sty are NOT in this pipeline's
%% TeX Live image (only the sibling ijdc-v9/idcc classes were checked and
%% also absent), so both are bundled here, extracted from the upstream
%% dccpaper.dtx via docstrip in the pinned compiler itself -- byte-identical
%% to what a full TeX Live install would produce.
%%
%% dccpaper-by.pdf is ALSO bundled: the class's own \maketitle unconditionally
%% stamps a Creative Commons Attribution 4.0 badge and notice onto every
%% article's title page (it licenses the OUTPUT ARTICLE, not the class
%% itself -- the class text is LPPL-1.3c). Without the image file the
%% \includegraphics call in dccpaper-base.sty fails.
%%
%% BIBLIOGRAPHY: dccpaper-base.sty supports both biblatex and apacite. This
%% starter uses apacite (classic bibtex-based) rather than biblatex, because
%% this pipeline's compile loop runs bibtex, not biber.
%%
%% Class documentation: https://ctan.org/pkg/dccpaper
%% Licence: dccpaper is LPPL-1.3c, (c) Digital Curation Centre, University of
%% Edinburgh. The bundled dccpaper-by.pdf is Creative Commons' own badge
%% image, used here exactly as the class intends: to mark the OUTPUT
%% article's licence, not to relicense this template.
\documentclass[conference]{ijdc-v14}
\usepackage[bibnewpage,numberedbib]{apacite}

\title{Preserving Reproducibility in Long-Running Simulation Studies}
\author{Ada Lovelace}
\affil{Digital Curation Institute, Institute of Analytical Engines}
\author{Alan Turing}
\affil{School of Computing Machinery}
\correspondence{Ada Lovelace, Digital Curation Institute, 1 Analytical
  Engines Walk, Cambridge, UK. Email: \email{ada@example.edu}}

\conference*{the International Digital Curation Conference IDCC26,
  16-18 February 2026}
\submitted{1 September 2025}

\begin{document}
\maketitle

\begin{abstract}
Replace this with your own abstract. State the curation problem you address,
what you did, and what you found.
\end{abstract}

\section{Introduction}
Start here. Cite things like this \cite{knuth1984texbook}.

\section{Related Work}
Summarise prior curation work relevant to your study.

\section{Method}
Describe your curation and preservation workflow.

\begin{table}[htbp]
  \centering
  \caption{Example dataset inventory.}
  \label{tab:example}
  \begin{tabular}{lrr}
    \hline
    Dataset & Files & Size (GB) \\
    \hline
    Run A & 1,204 & 18.3 \\
    Run B & 3,981 & 62.7 \\
    \hline
  \end{tabular}
\end{table}

Table~\ref{tab:example} shows the dataset inventory.

\section{Conclusions}
The text of the conclusions starts here.

\section{Acknowledgements}
Any acknowledgements should be placed here.

\bibliographystyle{apacite}
\bibliography{references}

\end{document}
